\section{Modelo Matemático}
\label{sec:modelo}

% ==== REEMPLAZA por el modelo físico de TU caso (ver docs/casos/caso_tXX.md) ====
% Escribe la(s) ley(es) física(s) que generan el sistema (LVK, LCK, equilibrio de
% fuerzas, balance de materia, diferencias finitas, etc.) y muestra la matriz
% \mathbf{A} y el vector \mathbf{b} con los valores numéricos de tu caso.

\subsection{Ecuación del sistema}
El balance en el caso de estudio conduce al sistema $\mathbf{A}\mathbf{x}=\mathbf{b}$:

\begin{equation}
\begin{bmatrix}
a_{11} & a_{12} & \cdots & a_{1n} \\
a_{21} & a_{22} & \cdots & a_{2n} \\
\vdots & \vdots & \ddots & \vdots \\
a_{n1} & a_{n2} & \cdots & a_{nn}
\end{bmatrix}
\begin{bmatrix} x_1 \\ x_2 \\ \vdots \\ x_n \end{bmatrix}
=
\begin{bmatrix} b_1 \\ b_2 \\ \vdots \\ b_n \end{bmatrix}
\label{eq:matriz}
\end{equation}

% Sustituye los coeficientes por los valores numéricos de tu caso.

\subsection{Propiedades de la matriz}
\begin{itemize}
	\item \textbf{Tipo de matriz:} (simétrica / tridiagonal / diagonal dominante / densa).
	\item \textbf{Número de condición:} $\mathrm{Cond}(\mathbf{A}) = \|\mathbf{A}\|\,\|\mathbf{A}^{-1}\|$.
	\item \textbf{Solver recomendado:} (según la ficha técnica de tu caso).
\end{itemize}
